\documentclass[11pt]{article}
\usepackage[margin=0.82in]{geometry}
\usepackage{amsmath,amssymb,amsthm,graphicx,microtype}
\usepackage[hidelinks]{hyperref}

\title{Uniformly bounded translations on de Branges spaces:\\
a complete Paley--Wiener dichotomy}
\author{Full result for arXiv:1801.01345}
\date{13 August 2026}

\newtheorem{theorem}{Theorem}
\newtheorem{lemma}{Lemma}
\newcommand{\C}{\mathbb C}
\newcommand{\R}{\mathbb R}
\newcommand{\Hh}{\mathcal H}
\newcommand{\PW}{PW}

\begin{document}
\maketitle

\begin{center}
\fbox{\parbox{0.92\textwidth}{\textbf{Status: full classification, likely valid.}
Every nonzero de Branges space whose full real-translation group is uniformly
bounded is either a Paley--Wiener space as a set or the one-dimensional space
of constants.  Thus the intended infinite-dimensional converse is true; the
constants are the unique exception if the source question is read literally
without a dimension hypothesis.}}
\end{center}

\section*{1. Source question and answer}

Baranov and Bommier-Hato prove that a de Branges space has a representing
Fock weight depending only on the imaginary part exactly when it is a
Paley--Wiener space $\PW_a$, $a>0$.  They observe that this property implies
uniform boundedness of all real translations
\[
 U_tF(z)=F(z+t),\qquad t\in\R,
\]
and ask whether the converse holds \cite{source}.

\begin{figure}[h]
\centering
\includegraphics[width=0.96\textwidth]{figures/open_problem_crop.png}
\caption{Theorem 2.4 and the ensuing open question on page 5 of
arXiv:1801.01345.}
\end{figure}

The following theorem completely answers the question, including its
finite-dimensional edge case.

\begin{theorem}[Uniform-translation dichotomy]\label{thm:main}
Let $\Hh$ be a nonzero de Branges space and suppose that every $U_t$ maps
$\Hh$ onto itself and
\[
 C:=\sup_{t\in\R}\|U_t:\Hh\to\Hh\|<\infty.
\]
Then exactly one of the following holds:
\begin{enumerate}
\item $\Hh=\PW_a$ as a set, for some $a>0$, and its norm is equivalent to
the usual Paley--Wiener norm;
\item $\Hh=\C$ as a set.
\end{enumerate}
In particular, every infinite-dimensional such space is a Paley--Wiener
space.  The second case is realized by $\Hh(z+i)$ and is the unique literal
exception to the converse of Theorem 2.4 in the source.
\end{theorem}

\section*{2. Proof intuition}

Uniform boundedness can first be removed by averaging the original inner
product over longer and longer translation intervals.  The resulting
equivalent norm makes every translation unitary.  The decisive point is that
this averaging preserves all de Branges axioms: conjugation commutes with
translations, while a zero-moving Blaschke factor at $w$ is carried by a
translation to the corresponding factor at $w-t$.

For the unitary norm, the reproducing kernel is invariant under simultaneous
real translation of both variables.  If $E=A-iB$ generates the space, the row
$V=(A,B)$ must therefore satisfy
\[
 V(z+t)=V(z)M_t,qquad M_t\in SL(2,\R).
\]
The matrices form a continuous one-parameter group, hence $M_t=e^{tX}$ for
$X\in\mathfrak{sl}(2,\R)$.  Cayley--Hamilton leaves elliptic, parabolic, and
hyperbolic possibilities.  They give respectively a sinc kernel, a constant
kernel, and a hyperbolic-sinc kernel.  The first is Paley--Wiener, the second
is one-dimensional, and the third violates the two-point Cauchy--Schwarz
inequality for positive kernels.

\section*{3. Averaging lemma}

Write $J F=F^\#$ and, when $F(w)=0$ for $w\notin\R$, put
\[
 R_wF(z)=\frac{z-\bar w}{z-w}F(z).
\]

\begin{lemma}\label{lem:average}
Under the hypotheses of Theorem~\ref{thm:main}, there is an equivalent Hilbert
norm $\|\cdot\|_0$ on the same functions such that:
\begin{enumerate}
\item every $U_t$ is unitary for $\|\cdot\|_0$;
\item $J$ is an isometry for $\|\cdot\|_0$;
\item $\|R_wF\|_0=\|F\|_0$ whenever $F(w)=0$ and $w\notin\R$.
\end{enumerate}
Consequently the new norm is again a de Branges norm.
\end{lemma}

\begin{proof}
Since $U_t^{-1}=U_{-t}$,
\begin{equation}\label{eq:twosided}
 C^{-1}\|F\|\leq\|U_tF\|\leq C\|F\|.
\end{equation}
The scalar coefficient functions of $U_t$ are continuous.  Indeed this is
clear when the second vector is a reproducing kernel, because
$\langle U_tF,K_w\rangle=F(w+t)$; density of the kernel span and the uniform
bound extend continuity to every second vector.

For $T>0$, define the bounded positive operator $A_T$ weakly by
\[
 \langle A_TF,G\rangle=
 \frac1{2T}\int_{-T}^{T}\langle U_tF,U_tG\rangle\,dt.
\]
Equation \eqref{eq:twosided} gives
\[
 C^{-2}I\leq A_T\leq C^2I.
\]
Choose a weak-operator cluster point $A$ as $T\to\infty$.  Translating the
averaging interval by a fixed $s$ changes only two boundary intervals of
total length $2|s|$; the integrand is uniformly bounded.  Hence
\[
 U_s^*AU_s=A\qquad(s\in\R).
\]
Moreover $C^{-2}I\leq A\leq C^2I$.  Thus
\[
 [F,G]_0:=\langle AF,G\rangle,
 \qquad \|F\|_0^2=[F,F]_0,
\]
is an equivalent Hilbert norm and all $U_s$ are unitary.

The original de Branges conjugation is isometric and $JU_t=U_tJ$.  Taking
the same weak limit of averaged squared norms shows
$\|JF\|_0=\|F\|_0$.  Likewise,
\[
 U_tR_wF=R_{w-t}U_tF.
\]
When $F(w)=0$, the function $U_tF$ vanishes at $w-t$, so the original de
Branges axiom gives $\|R_{w-t}U_tF\|=\|U_tF\|$ for every $t$.  Averaging and
passing to the cluster point yields $\|R_wF\|_0=\|F\|_0$.
Point evaluations remain bounded because the norms are equivalent.  These are
precisely the de Branges axioms \cite{deBranges}.
\end{proof}

\section*{4. The $SL(2,\R)$ orbit}

By Lemma~\ref{lem:average}, represent the newly normed space isometrically as
$\Hh(E)$ for a Hermite--Biehler function $E=A-iB$, where $A^\#=A$ and
$B^\#=B$.  Set
\[
 V(z)=(A(z),B(z)).
\]
The de Branges reproducing-kernel formula is
\begin{equation}\label{eq:dbkernel}
 K(z,w)=-\frac{\det\!\begin{pmatrix}V(z)\\V(\bar w)\end{pmatrix}}
 {\pi(z-\bar w)}.
\end{equation}

Because translations are unitary,
$K(z+t,w+t)=K(z,w)$.  The denominator in \eqref{eq:dbkernel} is unchanged,
so all the determinants in its numerator are unchanged.  Choose real
$x_1,x_2$ such that $V(x_1),V(x_2)$ are linearly independent.  Such points
exist, since otherwise $A$ and $B$ would be proportional and the kernel would
vanish.  There is a unique real matrix $M_t$ with
\[
 V(x_j+t)=V(x_j)M_t\qquad(j=1,2).
\]
Preservation of their determinant gives $\det M_t=1$.  Comparing determinants
against the two independent translated rows then gives
\begin{equation}\label{eq:orbit}
 V(z+t)=V(z)M_t\qquad(z\in\C).
\end{equation}
The matrices depend continuously on $t$, and uniqueness in
\eqref{eq:orbit} gives $M_{s+t}=M_sM_t$.  Therefore
\[
 M_t=e^{tX}\quad\text{for some }X\in\mathfrak{sl}(2,\R).
\]
Putting $v=V(0)$ and using analytic continuation from real $z$ in
\eqref{eq:orbit},
\begin{equation}\label{eq:exporbit}
 V(z)=ve^{zX},\qquad z\in\C.
\end{equation}
The rows $v$ and $vX$ are independent; otherwise the orbit lies in one line
and again $K\equiv0$.

\section*{5. Classification and proof of the theorem}

Since $\operatorname{tr}X=0$, Cayley--Hamilton says
\[
 X^2=-(\det X)I.
\]
Write $e^{uX}=c(u)I+s(u)X$.  From \eqref{eq:exporbit}, multiplication of both
determinant rows in \eqref{eq:dbkernel} by $e^{-\bar wX}$ gives
\begin{equation}\label{eq:threeforms}
 K(z,w)=\gamma\frac{s(z-\bar w)}{z-\bar w},
 \qquad \gamma=K(0,0)>0,
\end{equation}
where $s'(0)=1$ and exactly one of the following occurs:
\[
\begin{array}{c|c}
\det X=a^2>0&s(u)=\sin(au)/a\\[2mm]
\det X=0&s(u)=u\\[2mm]
\det X=-a^2<0&s(u)=\sinh(au)/a.
\end{array}
\]

In the elliptic case, \eqref{eq:threeforms} is a positive scalar multiple of
the standard Paley--Wiener kernel
\[
 \frac{\sin(a(z-\bar w))}{\pi(z-\bar w)}.
\]
Hence the newly normed space, and therefore the original space as a set, is
$\PW_a$.

In the parabolic case the kernel is the constant $\gamma$, so the RKHS is
exactly the one-dimensional space of constants.  It is realized as the de
Branges space $\Hh(z+i)$.

The hyperbolic case is impossible.  For distinct real $x,y$,
\[
 K(x,x)=K(y,y)=\gamma,
 \qquad |K(x,y)|=gamma\frac{\sinh(a|x-y|)}{a|x-y|}>\gamma,
\]
contradicting $|K(x,y)|^2\leq K(x,x)K(y,y)$ for a positive kernel.
This proves Theorem~\ref{thm:main}.

\section*{6. Verification, novelty, and scope}

The proof was checked at four structural points: the two-sided bounds survive
the weak-operator limit; averaging preserves the zero-moving axiom because
$U_tR_w=R_{w-t}U_t$; determinant invariance produces a genuine continuous
$SL(2,\R)$ group; and the hyperbolic kernel has a negative two-point Gram
determinant.  The three matrix-exponential and kernel forms were also checked
symbolically by \texttt{code/sl2\_kernel\_audit.py}.  This code is a mechanical
audit only; no computational assertion is needed for the proof.

The classification distinguishes the source's two possible readings.  In the
intended infinite-dimensional setting, the missing implication is true and,
by the source's Theorem 2.4, a representing weight depending only on the
imaginary part exists.  Without a dimension condition, $\Hh(z+i)=\C$ is a
counterexample to that implication, but the theorem proves it is the only one.

The run indexes contained no duplicate.  Bounded searches through 13 August
2026 covered the exact question, its title and authors, translation-invariant
and periodic de Branges spaces, Bellavita's 2021 and 2023 translation-operator
papers, the 2024 homogeneous-space classification, and the 2026 composition-
operator paper arXiv:2604.06801.  The later works give criteria for individual
translations or other group actions, but no resolution of this uniformly
bounded full-group question was found.  Novelty outside this bounded search is
not asserted pending expert review.

\paragraph{Human-review recommendation.}
Verify the weak-operator averaging step, the preservation of the de Branges
zero axiom, and the elementary passage from determinant invariance to
$M_t\in SL(2,\R)$.  Also preserve the explicit one-dimensional caveat when
stating the answer to the source.

\begin{thebibliography}{9}
\bibitem{source}
A. Baranov and H. Bommier-Hato,
\emph{De Branges spaces and Fock spaces},
Complex Variables and Elliptic Equations 63 (2018), 907--930;
arXiv:1801.01345.

\bibitem{deBranges}
L. de Branges,
\emph{Hilbert Spaces of Entire Functions},
Prentice--Hall, Englewood Cliffs, 1968.
\end{thebibliography}

\end{document}
